\documentclass[a4paper,12pt]{ctexart}
\usepackage{xcolor}
\usepackage{graphicx}
\usepackage[top=1.8cm, bottom=1.8cm, left=1.8cm, right=1.8cm]{geometry}
\usepackage{array}
\usepackage{hyperref}
\usepackage{booktabs}
\hypersetup{colorlinks=true,linkcolor=blue!70!black}
\linespread{1.35}

\title{\textcolor{blue!70!black}{\Huge\textbf{计算机科学与技术}}}
\subtitle{\textcolor{gray}{信息技术与互联网方向 · 学生版}}
\author{}
\date{}

\begin{document}
\maketitle
\thispagestyle{empty}

\section{一句话概括}

计算机科学与技术研究怎么用算法和程序让计算机解决问题。从写App到开发操作系统，从训练AI模型到设计芯片架构，都是它的范围。

\section{自我评估——你适合吗？}

\begin{itemize}
  \item[\textbf{☐}] 碰到电脑问题第一反应是自己搜、自己试，不直接喊"坏了"
  \item[\textbf{☐}] 数学还可以，尤其逻辑和代数不觉得难
  \item[\textbf{☐}] 能对着一个debug问题坐两小时不崩
  \item[\textbf{☐}] 英语阅读没太大障碍（大量文档和社区都是英文）
  \item[\textbf{☐}] 喜欢搭积木式的创造感——从零搭出一个能跑的东西
  \item[\textbf{☐}] 对新工具、新框架有好奇心，愿意持续学
\end{itemize}

\textbf{勾了4个以上}：这个专业值得认真考虑。
\textbf{勾了2-3个}：可以试试，但要想清楚自己哪里需要补。
\textbf{勾了1个以下}：建议看看别的方向。

\section{大学四年学什么}

\begin{tabular}{ll}
\toprule
大一 & 程序设计基础（C/C++）、高等数学、线性代数、离散数学 \\
大二 & 数据结构、算法设计与分析、计算机组成原理、操作系统 \\
大三 & 计算机网络、数据库系统、软件工程、机器学习、编译原理 \\
大四 & 毕业设计、实习、前沿选修（AI安全、分布式系统等） \\
\bottomrule
\end{tabular}

\section{北京高考选科要求}

\begin{tabular}{ll}
\toprule
院校 & 选科要求 \\
\midrule
清华计算机 / 北大信科 & 物理必选 \\
北航计算机 / 北理工 & 物理必选 \\
北邮计算机 & 物理必选 \\
北交大 / 北工大 & 物理必选 \\
首师大 / 北联大 & 物理建议（部分不限） \\
\bottomrule
\end{tabular}

\section{这个专业最大的几个坑}

\begin{enumerate}
  \item 35岁焦虑是真的吗？——纯写代码的职业生涯确实有天花板，但转管理、架构、产品都是路
  \item 计算机不是"会装系统就行"——本科数学和理论基础课很硬，不是一直敲代码
  \item 自学能力决定上限——课堂教的永远滞后于工业界，最厉害的东西得自己课外学
  \item 竞争在加剧——十年前培训班出来就能找工作，现在大厂校招基本都是科班硕士起步
\end{enumerate}

\end{document}
